\section{Conclusion}

This paper evaluates weak-form EMH through benchmark-relative forecast testing, not structural market-efficiency identification. Under expanding-window evaluation, leakage controls, multiple loss metrics, and a stronger benchmark family, random walk remains the best-performing model by RMSE across all tested \((L,H)\) settings.

Conditional diffusion does not provide incremental predictive gains in this implementation and is significantly worse than random walk in all Diebold--Mariano comparisons. Diffusion ablation and seed analyses also indicate sensitivity rather than stable superiority. XAI outputs, interpreted as model reliance diagnostics, are not stable enough to support strong economic narratives.

The main contribution is therefore methodological: a transparent, reproducible, benchmark-disciplined pipeline for weak-form predictability testing in Latin American ETF proxies, including explicit robustness boundaries and interpretation discipline. Stronger claims would require broader benchmark classes, deeper hyperparameter validation, and an explicit post-cost economic significance layer.
